\documentclass[11pt]{article}
\usepackage[margin=1in]{geometry}
\usepackage{amsmath,amssymb,amsthm,booktabs,enumitem,graphicx}
\newtheorem{theorem}{Theorem}\newtheorem{proposition}[theorem]{Proposition}
\newcommand{\fn}[1]{\texttt{\detokenize{#1}}}
\title{Lemma O, v3.2: certified thresholds, hand-off hypotheses, and the closure H6}\author{}\date{}
\begin{document}\maketitle
\noindent\textbf{v3.2 (12 Nov 2026).} Changes against v3 (\fn{../lemmaO/lemmaO_v3.tex}, kept intact): the cone column is evaluated at the needed ratio $1.6001\cdot10^{-7}$ (v3 printed deviations at the withdrawn $1.2\cdot10^{-7}$); the far family is $(M,K)=(0.15,0.3)$; H5 is replaced by data thresholds; H6 is proved (\S\ref{sec:H6}) conditionally on the strip statement of Lemma S(B), which is not proved (far-field localisation gap, \S\ref{sec:open}); Lemma~R (\fn{lemmaR.tex}) is new. Files: \fn{lemmaO_v32_interval.py v32} (python-flint \texttt{arb}, 200 bits), output \fn{lemmaO_v32_consts.json}; \fn{lemmaO_v32_h6.py}, output \fn{lemmaO_v32_h6.json}. Formulas: those of \texttt{lemmaO\_v2.tex} (constants $\Lambda_1,C_*,C_1,C_2,C_3$ on the cylinder $[-X,X]$, thresholds $t_j^*$) with one change.

\section{What changed}
\textbf{(a) Sharper $\kappa$.} Hypothesis (A) at $t=0$ gives $\|h_0'''\|_{L^2(-X,X)}\le K_3$, so $\|g_{xxx}\|_{L^2}\le\kappa:=K_3+\min\{\sqrt{2X}\,M,\,K_3\}$ (v2 used $K_3+\sqrt{2X}M$). Every other formula is unchanged; $\kappa$ enters monotonically, so the proof is unchanged.
\textbf{(b) Certification.} All constants are computed in ball arithmetic from exact rational inputs; the closed-form sups ($r=(16/17)^3$, $A_1,B_0,B_1,C_1',\dots,C_4'$, $S_1=15/8$, $S_2=10/\sqrt3$) are proved in v1 Lemma 1. The only non-closed-form step, $C_*=2\max_{u\ge0}f(u)$ with $f=-a_{\min}u^8+\sum_{k\le6}r_ku^k$, $r_k\ge0$, is bounded rigorously: (i) tail: if $\sum_kr_kU^{k-8}<a_{\min}$ then $f<0$ on $[U,\infty)$ (each $u^{k-8}$ decreases); (ii) on $[0,U]$ branch-and-bound with ball enclosures of $f$ on subintervals (relative tolerance $10^{-6}$); the returned value is the maximum of the leaf upper bounds. Each of $C_1,C_2$ is the minimum of two valid upper bounds. Pitfalls avoided: every radicand is a ball of a quantity that is a sum/product of positive exact data (no \texttt{ball**n} on balls straddling $0$, no \texttt{abs}); powers $1/2,1/4,3/2$ are \texttt{sqrt} compositions. $t_j^*$ are reported as \emph{lower} bounds (upper bounds on $C_j$ in the denominator). The script reproduces the double-precision v2 values to $10^{-6}$ relative (self-check in \texttt{\_\_main\_\_}).
\textbf{(c) Meaning of $\rho\le54,\ 67$.} The lemma is scale invariant, so $\rho$ does not enter; I used only $(X,M,K_2,K_3)$ and the time ratio. Please confirm $M=0.12$ is also the right value for $X=2.4$ (the sweep in the JSON gives $t^*$ for $M\in[0.08,0.3]$).

\section{Slack convention}
$K_2=K_3=(1+f)M$, $f=\frac12$. Closure of the continuity argument needs, for $t\le t_1$: $C_2\sqrt t\le fM$ (pointwise $h_{xx}$), $C_1\sqrt t\le \frac14-\sup|h_{0,x}|$ (slope), and for the $L^2$ bound on the \emph{whole} cylinder (the lemma controls $g_{xxx}$ on $[-1,1]$ only; cover $[-X,X]$ by $\lceil X\rceil$ intervals $[c-1,c+1]$): $\|h_{0,xxx}\|_{L^2(-X,X)}+\sqrt{\lceil X\rceil C_3t}\le K_3$. The tabulated $t^*=\min(t_1^*,t_2^*,t_3^*)$ uses the pointwise/density slack $t_j^*=(fM/C_j)^2$, $t_3^*=2(fM)^2/C_3$; the global $L^2$ requirement is listed separately below (it is a \emph{hypothesis on $h_0$}, not a time restriction).

\section{Near the S1 cone (certified)}
\begin{center}\small\begin{tabular}{lcc}\toprule
 & $X=3$ & $X=2.4$\\\midrule
$M$, $K_2=K_3$ & $0.12,\ 0.18$ & $0.12,\ 0.18$\\
$w$ (cutoff width) & $0.271$ & $0.194$\\
$\kappa$ & $0.36$ & $0.36$\\
$C_1\le$, $C_2\le$, $C_3\le$ & $18.70,\ 98.87,\ 3574$ & $26.09,\ 171.6,\ 1.261\cdot10^4$\\
$t_1^*\ge$ (slope) & $1.030\cdot10^{-5}$ & $5.29\cdot10^{-6}$\\
$t_2^*\ge$ ($h_{xx}$, binding) & $3.683\cdot10^{-7}$ & $1.223\cdot10^{-7}$\\
$t_3^*\ge$ ($g_{xxx}$ density) & $2.014\cdot10^{-6}$ & $5.71\cdot10^{-7}$\\
\textbf{certified} $t^*\ge$ & $\mathbf{3.683\cdot10^{-7}}$ & $\mathbf{1.223\cdot10^{-7}}$\\
needed, incl.\ $(1-4\bar a)^{-1}$ & $1.6001\cdot10^{-7}$ & $5\cdot10^{-8}$\\
\textbf{margin} (factor) & $\mathbf{2.30}$ & $\mathbf{2.44}$\\
at the needed time: $C_0\sqrt t$, $C_1\sqrt t$ & $0.00227,\ 0.00748$ & $0.0015,\ 0.0059$\\
 $C_2\sqrt t$ vs slack $0.06$ & $0.03955$ & $0.039$\\
 $\sqrt{C_3t}$ per conclusion region ($X=2.4$ column: $\sqrt{\lceil X\rceil C_3t}$ at $5\cdot10^{-8}$) & $0.02392$ & $0.0435$\\
slope: $\tfrac14-M-C_1\sqrt t\ \ge$ & $0.1225$ & $0.1241$\\\bottomrule
\end{tabular}\end{center}
The $X=3$ column is the v3.2 run ($w$ grid of 14 points; v3 had $w=0.317$, $t^*\ge3.701\cdot10^{-7}$, also valid). The $X=2.4$ column is unchanged from v3 and not used.

Sweep ($K=1.5M$, certified $t^*$): $X=3$: $3.93,3.82,3.68,3.44,3.03,2.36\cdot10^{-7}$ for $M=0.08,0.1,0.12,0.15,0.2,0.3$; $X=2.4$: $1.25,1.24,1.22,1.17,1.07,0.88\cdot10^{-7}$.

\section{Far-curve regime}
The slack-$\frac12M$ certificate at $K,M\le10$ is \textbf{not} available at time ratio $10^{-9}$: certified $t^*\ge1.2\cdot10^{-10}$ ($M=6.67,K=10$, $X=2$; $1.3\cdot10^{-10}$ at $X=2.4$), i.e.\ short by a factor $8.3$ (resp.\ $7.6$); $M=5,K=10$ ($f=1$): $3.0\cdot10^{-10}$, short by $3.4$; $M=10,K=15$: $1.8\cdot10^{-11}$, short by $56$. What \emph{is} certified at $10^{-9}$ (script, each re-certified at the stated rational $M$):
\begin{center}\small\begin{tabular}{lccc}\toprule
cylinder, $f=\frac12$ & $M\le$ & $K_2=K_3=1.5M$ & certified $t^*\ge$\\\midrule
$X=2$ ($w=0.19$) & $3.69$ & $5.535$ & $1.0011\cdot10^{-9}$\\
$X=2.4$ ($w=0.20$) & $4.10$ & $6.15$ & $1.0019\cdot10^{-9}$\\\bottomrule
\end{tabular}\end{center}
Monotonicity: all of $\Lambda_1,C_*,C_1,C_2,C_3$ are nondecreasing in $K_2,K_3,M$ (nonnegative coefficients, $\kappa$ increasing), so with $K$ fixed at the level above the certificate holds for every $M'\le M_{\max}$ (slack $K-M'\ge fM_{\max}$) at the same $w$. Deviations at $t=10^{-9}$ ($X=2$): $C_0\sqrt t\le0.033$, $C_1\sqrt t\le0.14$, $C_2\sqrt t\le1.84$ (slack $1.845$), $\sqrt{\lceil X\rceil C_3t}\le3.41$ so the global $L^2$ condition is $\|h_{0,xxx}\|_{L^2(-2,2)}\le2.12$ ($X=2.4$: $\le1.39$). Using $f=1$ leaves essentially no room for the global $L^2$ term ($0.005$ at $X=2$, negative at $X=2.4$), so $f=\frac12$ is the usable choice.

\section{Hypotheses the step-4 hand-off must check (per cylinder, rescaled so the cylinder is $[-X,X]$)}
\begin{enumerate}
\item[H1] \textbf{Graph.} After translating to the center and rotating so the tangent is horizontal (for the bent cone: the local tangent frame), $\gamma\cap([-X,X]\times[-H,H])$ is the graph of $h_0$ over $[-X,X]$ with $\sup|h_0|\le H-C_0\sqrt{t_1}$ ($C_0\sqrt{t_1}\le0.0023$ cone, $\le0.05$ far; v3.2 far value $9.3\cdot10^{-4}$).
\item[H2] \textbf{Slope.} $\sup_{[-X,X]}|h_{0,x}|\le\frac14-C_1\sqrt{t_1}$; v3.2: $|h_{0,x}|\le s_0=0.08$ in both families (Table~\ref{tab:thr}); this is below $M$ and below $\frac14-C_1\sqrt{t_1}$ ($0.2425$ cone, $0.2456$ far).
\item[H3] \textbf{Curvature jets.} $|h_0^{(j)}|\le M$ on $[-X,X]$, $j=2,3,4$ (and $j=1$ as in H2), with $M=0.12$ (cone) resp.\ $M=0.15$ (far, v3.2).
\item[H4] \textbf{Fifth derivative.} $\int_{-X}^X(h_0^{(5)})^2\le M^2$ (``good-time'' selection).
\item[H5] \textbf{Data thresholds (v3.2).} In own cylinder units on the whole cylinder: slope $\le s_0$, $\sup|h_0''|\le m_2$, $\|h_0'''\|_{L^2(-X,X)}\le n_3$ with the values of Table~\ref{tab:thr}. (v3: $\|h_{0,xxx}\|_{L^2}\le K_3-\sqrt{\lceil X\rceil C_3t_1}$, i.e.\ $0.14$ at $1.2\cdot10^{-7}$, $0.1337$ at $1.6\cdot10^{-7}$; superseded, since the closure is now Theorem~\ref{thm:H6}.)
\item[H6] \textbf{A priori bounds close.} Proved in \S\ref{sec:H6} (Theorem~\ref{thm:H6}) from H1--H5, H7, H9 and the strip statement of Lemma S(B) at levels $0$--$2$; no longer a hand-off check. See \S\ref{sec:open} for the one open input.
\item[H7] \textbf{Time.} $t_1\le$ needed ratio (cone: $1.6\cdot10^{-7}=50^{-4}$ at $X=3$ for $y_O=200$, half-width $0.75r$, i.e.\ $L=50\lambda_1$, $t_1=1$; $5\cdot10^{-8}$ at $X=2.4$; far: $L=113\lambda_1$, ratio $6.134\cdot10^{-9}$ with $M=0.15$, $K=0.3$, certified $t^*\ge9.511\cdot10^{-8}$, margin $15.5$), measured in the \emph{same} unit as $X$ (time ratio $=t/L^4$ for cylinder length unit $L$).
\item[H8] \textbf{Other arms} (not a requirement in step 3: the lemma is applied arc by arc to an immersed curve)\textbf{.} Distance (in cylinder units) from the other arm to the graph at time $0$ $\ge C_0\sqrt{t_1}+C_0'\sqrt{t_1}$, where $C_0'$ is the displacement bound of the cylinder centered on the other arm; $\ge0.1$ suffices in every case above ($C_0\sqrt{t_1}\le0.05$).
\item[H9] \textbf{Regularity.} $h$ smooth on $[-X,X]\times[0,t_1]$ (graph over the whole interval persists).
\end{enumerate}
\section{Revision (step-3 assembly)}
\textbf{Time ratio.} The ratio is $t_1/L^4$, $t_1=1$ ($\lambda_1^4$), $L=$ half-width$/X$. Cone, $y_O=200$, half-width $0.75r$, $X=3$: $L=50$, ratio $1.6\cdot10^{-7}$, margin $2.30$ ($1.2\cdot10^{-7}$ above corresponded to half-width $161\lambda_1$ and is withdrawn). \textbf{Far cylinders (v3.2).} $X=2$, $L=113\lambda_1$, $M=0.15$, $K_2=K_3=0.3$, $w=0.136$: $C_1,C_2,C_3\le55.83,\ 486.4,\ 1.350\cdot10^5$, certified $t^*\ge9.511\cdot10^{-8}$ against $6.134\cdot10^{-9}$ (margin $15.5$); $C_0\sqrt t\le9.3\cdot10^{-4}$, $C_1\sqrt t\le0.00438$, $C_2\sqrt t\le0.03810$, $\sqrt{C_3t}\le0.02878$. Physical requirement ($\lambda_1$ units): $|h_{xx}|\le3.53\cdot10^{-4}$, $|h_{xxx}|\le1.17\cdot10^{-5}$, $|h_{xxxx}|\le1.03\cdot10^{-7}$, $\|h_{xxx}\|_{L^2(|x|\le2L)}\le3.32\cdot10^{-6}$, slope $\le0.08$, $\int(h_0^{(5)})^2\le0.0225$ (cylinder units). The v3 far case $(M,K)=(1,1.5)$ (recomputed: $t^*\ge1.249\cdot10^{-8}$, margin $2.04$, $C_2\sqrt t\le0.3503$) is withdrawn: with it H6 fails ($\times0.54$ for the cone $K_2$). \textbf{$K_3$ hand-off.} Cone: $|h_0^{(j)}|\le0.12$, the thresholds of Table~\ref{tab:thr}, a priori $K=0.18$; Lemma S, (B) windows with the bent-cone profile (all energies are of the deviation $v=f-U_e$ from the profile frozen at entry, see Lemma S v3.2; $K$'s include the $U$-part): $E_{1,2,3}\le3\cdot10^{-6}$ on windows $y\in[50,253]$ and $E_4\le3\cdot10^{-6}$ at $\tau_1$ give $K_2\le0.044$, $K_3\le0.014$, $K_4\le0.041$ in cylinder units $L=50$ (margins $2.7,8.4,2.9$ against $M=0.12$; a priori $K=0.18$ has $\ge4.1$). The previous levels $(10^{-4},3\cdot10^{-4},5\cdot10^{-4})$ and $3\cdot10^{-5}$ are withdrawn ($3\cdot10^{-5}$ gives $K_4=0.127>M$). Data: \fn{lemmaO_v32_consts.json}, keys \fn{cases}, \fn{far_v32}, \fn{far_v31_reference}.
\section{v3.2: data thresholds and the closure H6}\label{sec:H6}
Units $\lambda_1$; $r$ is the distance along the $S_1$ ray. Cone members: centres $r_m\in[200,380]$, unit $L_m=r_m/4$, $X=3$, cylinder $[r_m/4,7r_m/4]$, conclusion region $[3r_m/4,5r_m/4]$. Far members: centres $\ge376$, unit $113$, $X=2$; their conclusion regions cover $r\ge263$. Strip: $r\in[50,150)$. Since $C_j$ do not depend on $t$, a cone member of unit $L_m$ has the deviations of \S3 times $(50/L_m)^2$.
\begin{table}[h]\centering\small
\begin{tabular}{lccc}\toprule
 & slope $s_0$ & $\sup|h_0''|\le m_2$ & $\|h_0'''\|_{L^2}\le n_3$\\\midrule
cone & $0.08$ (data $0.0573$, $\times1.40$) & $0.03$ ($0.0189$, $\times1.59$) & $0.004$ ($1.20\cdot10^{-3}$, $\times3.3$)\\
far & $0.08$ ($0.0428$, $\times1.87$) & $0.04$ ($0.0208$, $\times1.93$) & $0.004$ ($7.2\cdot10^{-4}$, $\times5.6$)\\\bottomrule
\end{tabular}\caption{New data thresholds; in brackets the step-4 certified value and its margin.}\label{tab:thr}\end{table}

\noindent The thresholds are needed: with the data bounded only by $M=0.12$ the closure fails ($0.12+1.9\cdot0.0395>0.18$).

\begin{theorem}[H6]\label{thm:H6}
Assume (H5) as amended, Proposition B$'$ of Lemma S v3.2 (\fn{step3_proof_v32.tex} \S3.7) on the strip windows at levels $0$--$2$ with the entry energies of (H3), and that the solution is smooth on $[t_1,t]$. Then for every cone and far cylinder, on its whole base and on $[t_1,t]$: slope $\le0.1354$ (cone), $0.0876$ (far); $|h_{xx}|\le0.1067$ (cone), $0.1313$ (far); $\|h_{xxx}\|_{L^2}\le0.0938$ (cone), $0.1255$ (far). In particular hypothesis (A) of Lemma O ($\tfrac14$, $K_2$, $K_3$) holds strictly in every cylinder.
\end{theorem}
\begin{proof}
Let $\bar t$ be the supremum of times up to which (A) holds in every cylinder. On $[t_1,\bar t]$:

\emph{1. Lemma O in every member} gives, on its conclusion region, $|g_x|\le C_1\sqrt t$, $|g_{xx}|\le C_2\sqrt t$, $\int_{-1}^1g_{xxx}^2\le C_3t$, $g=h-h_0$.

\emph{2. Frames.} With curvature $k$, tangent angle $\phi$ in a frame and $|\tan\phi|\le s$: $h_{xx}=\sec^3\!\phi\,k$, $h_{xxx}=\sec^4\!\phi\,(k_\sigma+3\tan\phi\,k^2)$, $dx=\cos\phi\,d\sigma$. Hence for a target frame $c$ and a source frame $m$
\[|h^c_{xx}|\le(1+s_c^2)^{3/2}|k|,\quad|k|\le|h^m_{xx}|,\quad|h^c_{xxx}|\le(1+s_c^2)^2\bigl(|h^m_{xxx}|+3(s_c+s_m)k^2\bigr),\quad dx_c\le(1+s_m^2)^{1/2}dx_m,\]
and from the ray chart $|k|\le|f'|$, $|k_\sigma|\le|f''|+3|f|f'^2$.

\emph{3. Covering.} A point $p$ of a cone cylinder $c$ lies in the strip ($r_p<150$), or in the conclusion region of the cone member $r_m=\min(4r_p/3,380)$ ($150\le r_p\le475$; $L_m\ge50$), or in that of the far member centred at $p$ ($r_p>475$). A point of a far cylinder lies in a cone conclusion region ($150\le r_p<263$) or in a far one.

\emph{4. Slope.} At a covered point the tangent angle changes by at most $C_1\sqrt t$, so the slope in any frame is $\le\tan(\arctan s_0+C_1\sqrt t)\le0.0876$. On the strip the angle to the ray is $\le\arctan\sup|f|$ and the axis of $c$ makes the angle $\arctan|f(r_c,\tau_1)|$ with the ray: slope $\le\tan(\arctan0.0420+\arctan0.0928)\le0.1354$.

\emph{5. $K_2$.} Covered: $L_c|k(p,t)|\le(1+s_0^2)^{3/2}m_2+(L_c/L_m)(50/L_m)^2C_2\sqrt t$; the data term carries no ratio because it is $L_c|k_0(p)|$. The worst case is $r_c=380$, $r_p=150$ ($L_m=50$, ratio $1.9$); beyond $475$ the deviation is $(L_c/113)C_2^{\rm far}\sqrt t$. Strip: $L_c|k|\le L_c\,y|f'|/y\le0.0497$.

\emph{6. $K_3$.} Split the base into the strip and the pieces $[150,250]$, $[250,416.7]$, $[416.7,475]$ (members $200$, $333.3$, $380$) and one far piece. A member piece $P$ contributes
\[(1+s_c^2)^2\Bigl[(1+s_m^2)^{1/4}(L_c/L_m)^{3/2}\bigl(n_3+\sqrt{C_3t}\,(50/L_m)^2\bigr)+3(s_c+s_m)(L_ck)^2|P|^{1/2}\Bigr],\]
and the strip, window by window, $L_c^{3/2}$ times $y_0^{-3/2}(\sqrt{E^u_2}+\sqrt{E^u_1})+\bigl(p(p-1)U+3(|f|+s_c)(yf')^2\bigr)\bigl(\int y^{-4}\bigr)^{1/2}$; the squares are added. Only levels $\le2$ of (B) enter.

\emph{7.} The bounds of the statement are strictly below $\tfrac14$, $K_2$, $K_3$ and uniform over the (compact) families, so $\bar t=t$.
\end{proof}
\begin{center}\small\begin{tabular}{lccc}\toprule
closure item & bound & against & margin\\\midrule
cone slope (strip worst; covered $0.0876$) & $0.1354$ & $0.25$ & $\times1.85$\\
far slope & $0.0876$ & $0.25$ & $\times2.86$\\
\textbf{cone $K_2$} (covered; strip $0.0497$) & $0.1067$ & $0.18$ & $\mathbf{\times1.69}$\\
cone $K_3$ (strip $0.0194$; pieces $0.0805$, $0.0235$, $0.0155$, $0.0340$) & $0.0938$ & $0.18$ & $\times1.92$\\
far $K_2$ & $0.1313$ & $0.3$ & $\times2.29$\\
far $K_3$ (pieces $0.1050$, $0.0233$, $0.0459$, $0.0459$) & $0.1255$ & $0.3$ & $\times2.39$\\
cone H4 (unchanged) / far H4 & $0.304M^2$ / $0.229M^2$ & $1$ & $\times3.28$ / $\times4.37$\\\bottomrule
\end{tabular}\end{center}
\emph{Robustness.} With all deviations taken at $\mu$ times the horizon every item still closes at $\mu=3.5$ (cone $K_2$, $K_3$ $\times1.04$), so $t^*$ ($\times2.30$) stays the binding time margin. With the strip energies multiplied by $2$, $4$, $10$, $30$ the smallest margin is $1.69$, $1.69$, $1.56$, $1.26$.

\noindent\emph{Order in the bootstrap} (D4/D5 of \fn{step3_proof_v32.tex}): D5a Lemma O conclusions from B4 and (H5); D4 windows, using D5a at their outer edge; D5b Theorem~\ref{thm:H6}, which improves B4. B4 reads: $K_2=K_3=0.18$ (cone), $0.3$ (far, $M=0.15$).

\section{Open input}\label{sec:open}
Theorem~\ref{thm:H6} uses the strip statement of Lemma S(B) at levels $0$--$2$ for all $t\in[t_1,T)$. Its outer cutoff ramp is bounded in v3.2 by Lemma~R (\fn{lemmaR.tex}; ramp $[190,379.5]\lambda_1$; certified contribution $1.6\%$, $0.09\%$, $3.05\times$ the entry energies at levels $0,1,2$, so the strip energies are at most $8$ times the per-window values and the margins of the table become $1.60$ (cone slope), $1.69$, $1.69$, $2.29$, $2.39$). \textbf{The localisation of the far-field energies is open} at every window edge in $60\lambda(\tau)\lesssim r<150\lambda_1$ (levels 1, 2; \fn{lemmaR.tex} \S5, \fn{../step3_gap_v32/gap_v32.tex}): the strip statement is not proved, in v3.1 or v3.2, and Theorem~\ref{thm:H6} is conditional on it. The same gap affects the crude window bounds beyond $y=60$ and the outer end of Lemma~S(A); Lemma~O and Lemma~R themselves are not affected. The theorem tolerates a strip energy $30$ times the entry value (smallest margin $1.26$).
\end{document}
